Where does the novelty a language model can produce with no task come from?
We test the three places usually credited, the sampler, the prompt and the
loop, on base models, with an LLM judge that is itself measured. An
entropy-banded anti-probable decoder roughly doubles $n$-gram novelty against
the training corpus but shows no detectable effect on judged surprise or
connection; input improbability shows no benefit under the
operationalizations we tested. The loop is different. Forced to continue a
premise with the end-of-text token masked, a base model degenerates into
repetitive modes; a closed loop built on the architecture of spontaneous
cognition revives it, and taking that loop apart over 23 conditions locates
most of the effect in two simple operations: a windowed repetition penalty
(habituation) and a periodic \emph{interruption}, one neutral sentence
injected every few hundred tokens. On windows of generated text only, judged
by Claude Opus~5 ($k=5$) with the premise as the unit ($n=10$), the
interruption raises surprise from 1.6 to 3.0 and connection from 1.3 to 3.7
over habituation alone (paired permutation $p=0.002$) and matches the full
scaffold. Part of that is self-copy: given the same four sentences in
rotation, the model replays its own earlier segments, out of the judge's
view; on fresh windows the interruption adds +1.2 to +1.4 surprise and
+0.8 connection, and a reset context, which cannot replay, adds +1.9 and
+2.0. Controls show what carries the effect: a bare paragraph break does
nothing and a continuity connective hurts; the new subject carries most of
it; injecting the premise or the stream's own past is as bad as not
interrupting; salience timing is no better than a clock. The gains are local:
judged as whole documents, no arm builds an integrated text. A pre-registered replication on ten new premises confirms
the primary contrast ($p=0.003$) and refutes the expectation that the kept
context carries connection. The ordering replicates
on three base models from two families, under a second judge family and in
the rankings of independent human readers. These results characterize a
simple, controllable intervention for forced open-ended generation; they do
not establish a general mechanism of creativity.
